% vim: set tw=72:
%!TEX root=./paper-dispersive-short-All.tex
\title{Testing quantum Darwinism dependence on observers' resources}

\author{A. Feller$^3$}
\author{B. Roussel$^2$}
\author{A. Pontlevy$^1$}
\author{P. Degiovanni$^1$}

\affiliation{(1) Univ Lyon, Ens de Lyon, Universit\'e Claude Bernard 
Lyon 1, CNRS, Laboratoire de Physique, F-69342 Lyon, France}
\affiliation{(2) Department of Applied Physics, Aalto University, Aalto, Finland}
\affiliation{(3) Univ. Lille, CNRS, Centrale Lille, UMR 9189 CRIStAL, F-59000 Lille, France}

\begin{abstract}
The emergence of an objective classical picture
is the core question of quantum Darwinism. How does
this reconstructed classical picture depends on the resources
available to observers?
%
In this Letter, we develop an experimentally relevant model of a qubit 
coupled dispersively to a transmission line and use time-frequency
signal processing techniques to understand if and how the emergent classical
picture is changed when we have the freedom to choose the
fragment decomposition and the type of radiation sent to probe
the system. We show the crucial role of correlations in the reconstruction 
procedure and point to the importance of studying the type of 
measurements that must be done to access an objective classical data.

\end{abstract}

\keywords{quantum Darwinism, decoherence, quantum optics, circuit QED,
time-frequency analysis}

\maketitle

%--------------------------------------------
% Introduction
%-------------------------------------------
%
Our view of the quantum-to-classical transition, originally based on
decoherence~\cite{Zurek-2003-1} has been expanded with the idea of quantum
Darwinism~\cite{Zurek-2009-1}. This new paradigm shifts the focus from 
the system $S$ to the observers that are probing it. The environment 
$E$ is then viewed as a noisy communication channel from which 
the information about the system is broadcasted to a network of 
observers. Despite being inherently many-body, the approach has led 
to many theoretical results~\cite{Zurek-2001,Riedel-2010,Riedel-2011-1,
Riedel-2012,Zurek-2013,Brandao-2015,Knott-2018-1,Zwolak-2014}
as well as experimental tests~\cite{Burke-2010-1,Zurek-2019}. Although
the existence of an objective observable has been shown to follow from
the first principles of quantum theory \cite{Brandao-2015,Knott-2018-1},
understanding if and how a classical image can be reconstructed by the
observers is still an open question
\cite{Ollivier-2022-1,Garcia-Pinto-2021-1}.

The central tenet of quantum Darwinism is to relate the emergence of a
consensus about the classical state of a quantum system to the many-body 
entanglement structure between the observers and the system. 
Zurek~\cite{Ollivier-2004-1,Zurek-2009-1} used the quantum mutual
information as an estimation of the accessible information \cite{Book-Wilde} in
asymptotic quantum Shannon theory \cite{Touil-2022-1} and
identified the existence of a plateau between the system and 
fragments of the environment as a signature of the redundant 
broadcast of information on pointer states, an intuition confirmed 
in~\cite{Le-2019-1,Feller-2021}.

However, the ability to reconstruct a classical image does not only
depends on the interaction between the system and the environment.
First, it also depends on the initial state of the environment, before
its interaction with the system, as demonstrated by recent studies on 
the role of thermalisation~\cite{Le-2021-1} and many-body 
localization~\cite{Mirkin-2021-1}. Second, it also depends on
the choice of a reference frame, that is which degrees of freedom are
accessible to which observer. The two reflect the resources available 
to observers to probe a quantum system on which we could add 
their computational and communication abilities. The foundational question of
the emergence of classicality can then be precised by analyzing to what 
extent it depends on these resources.

In this Letter, we study this resource dependence of quantum Darwinism. 
We consider a concrete model of a qubit dispersively coupled to a 
transmission line (see \cref{fig/the-idea}).
%and illuminated by a probe radiation (see \cref{fig/the-idea}). 
By using different quantum states for the probe radiation and different
time-frequency windows, 
both the influence of the reference frame and of the probe can be studied. 
We unravel the crucial role played by correlations in the classical
reconstruction process, mostly in a mesoscopic regime where the probe 
involves 1 to 10 photons on average. 

%------------------------------------------
% The dispersive model and time-frequency analysis
%------------------------------------------

\begin{figure}
\centering
\includegraphics{figures/scattering}
\caption{\label{fig/the-idea} 
A qubit $S$ is dispersively coupled to a transmission  
line $E$ such that incoming bosonic excitations at given energy 
$\hbar\omega$ are scattered by the qubit in the quantum state 
$\ket{s}$ via the scattering phase $\me^{\mi\theta_s(\omega)}$. 
%
Fragments are defined from a time-frequency atomic decomposition 
from which information about $\ket{s}$ is recovered.}
\end{figure}

\paragraph*{The dispersive model} The system $S$ is a qubit coupled
dispersively to a transmission line, 
acting as its environment $E$. 
%Their Hilbert spaces are respectively 
%denoted by $\mathcal{H}_S\approx \mathbb{C}^2$ and $\mathcal{H}_E$ 
%which is a bosonic Fock space. 
The transmission line is modeled by a 
chiral bosonic field with one chiral mode per pulsation $\omega>0$.
In the dispersive regime, coupling to the qubit is
described by an input/output relation 
$b_{\text{out}}(\omega)=\me^{\mi\theta_s(\omega)}\,b_{\text{in}}(\omega)$ 
where $\theta_s(\omega)$ is a pure phase~\footnote{As expected from
energy conservation in such a 1D chiral system.}.
This is realized in the low-power dispersive regime of circuit-QED 
\cite{Blais-2021-1} used for photon number 
\cite{Gambetta-2006-1,Schuster-2007-1} and qubit state measurements
\cite{Gambetta-2008-1}. 
For simplicity, we consider a time delay model with 
$\theta_s(\omega)=\omega\tau_s$ so that
%
\begin{equation}
b_{\text{out}}(\omega)= \me^{\mi \omega\tau_s} 
	\,b_{\text{in}}(\omega)
\,.
\label{eq/scattering}
\end{equation}
The two different states of the qubit $\ket{s}$ ($s=0,\ 1$) thus define a time delay 
difference $\Delta\tau=\tau_1-\tau_0$.
%

\emph{Fragments from atoms of signals}. Studying the emergence
of quantum Darwinism requires defining fragments of the environment 
$E$. We achieve this \emph{via} time-frequency atomic decompositions.
Concretely, an orthonormal basis $\ket{\phi_\alpha}$ of 
$L^2(\mathbb{R})$, $\alpha$ labeling a partition of the time-frequency plane,
defines a quantum reference frame. It corresponds to a 
mode decomposition for the quantum field in which each mode is an atomic 
fragment. Observers are defined 
by aggregating several atomic fragments. 
This enables us to change 
the atomic decomposition, \textit{i.e.} to change the reference frame.

For simplicity, we use the normalized Shannon atoms of signals defined 
as 
$w_{k,l}(\omega)=  \mathbf{1}_{\left[ \frac{2\pi}{T}k,
	\frac{2\pi}{T} (k+1) \right[}(\omega) \, \sqrt{\frac{T}{2\pi}}
	 \,\me^{\mi \omega lT}
\,,$
(see \cref{sec/notations} for normalizations)
with $(l,k)\in\mathbb{Z}\times \mathbb{N}$ providing a regular tessellation of 
the time-frequency half-plane in cells with a bandwidth 
$2\pi/T$ around $\omega_k=2\pi(k+1/2) / T$ and localized in time around
$t_l=lT$. The corresponding
annihilation and creation operators 
$b^{(\dagger)}_{k,l}$, are obtained from the $b^{(\dagger)}_\omega$
operators by a linear
unitary transformation. Since we are dealing with a bosonic
environment, the Fock space factorizes in a tensor product
$
\mathcal{H}_E = \bigotimes_{\omega\geq 0} \mathcal{H}_\omega
	=\bigotimes_{(k,l)\in\mathbb{N}\times\mathbb{Z}} \mathcal{H}_{k,l}
$
of the Fock spaces $\mathcal{H}_{k,l}$ associated with each atomic mode.

The ability to reconstruct a classical image of the system $S$ by
many independent observers in repeated experiments
is quantified by the accessible information
$I_{\text{acc}}(S,F)$ between $S$ and a fragment $F$
\cite{Touil-2022-1}. Apart from few cases \cite{Yao-2014-1}, the
accessible information involves a difficult optimization problem. 
Therefore, we will use the quantum mutual information 
$I(S,F) = S(S) + S(F) - S(SF)$ with $S(X)$ the quantum 
entropy of $X$ as a surrogate estimator \cite{Ollivier-2004-1}. 
%When the state
%of $SE$ is pure, we have $S(SF) = S(\overline{F})$ with $\overline{F}$
%the complementary set of $F$ in $E$. 
%The mutual information is our
%main quantity of interest.

%----------------------------------------
% Discussions of the results 
%---------------------------------------

\emph{Results} We considered classical and quantum probes.
The classical probe is a multi-mode coherent state
$
\ket{\phi} = \bigotimes_{\omega\geq 0} \ket{q\phi(\omega)}
	=\bigotimes_{(k,l)\in\mathbb{N}\times\mathbb{Z}} \ket{q\phi_{k,l}}
$
with wave-function $\phi_{k,l} = \int_0^{+\infty} 
w_{k,l}^*(\omega) \phi(\omega) \, \md \omega$ and
average photon number $\overline{n}=q^2$. This corresponds to the
illumination of an object with coherent radiation, a pure-state version
of Refs. \cite{Riedel-2010-1,Riedel-2011-1}. Because factorization
of a multi-mode coherent state holds in any reference frame, the
scattered radiation by a given state $\ket{s}$ of $S$ shows
no correlations between the various fragments.
%
The quantum probe is a Fock state $\ket{n[\phi(\omega)]}
= \frac{\left(a^\dagger[\phi] \right)^n}{\sqrt{n!}}\ket{0}$
with $n$ photons in the $\phi(\omega)$ mode. In contrast to 
coherent states, such a state is highly entangled in terms
of the atomic modes. 
%it is a superposition of many products
%of Fock states relative to them.
In both cases, the dispersive interaction \cref{eq/scattering} 
amounts to replace the incoming wavefunction by 
$\phi^{(s)}(\omega)=\me^{\mi\omega\tau_s}\phi(\omega)$
conditioned on the state $s$ of the system. 
The signal $\phi(\omega)$ will be conveniently represented
both in the time-domain 
$\phi(t) = \int_{\mathbb{R}} \phi(\omega) 
	\,\me^{\mi \omega t}
	\, \frac{\md \omega}{\sqrt{2\pi}}
$
but also in the time-frequency domain via its Wigner
representation 
$W_\phi(t,\omega) = \int_{\mathbb{R}} \phi^*(t-\tau/2)
\phi(t+\tau/2) \, \me^{\mi\omega\tau} \, \md\tau$.
%
For definiteness in numerical applications, 
a normalized Gaussian wavepacket
$\phi(\omega) = \frac{1}{\sqrt{\mathcal{N}}} 
	\,\me^{-\frac{(\omega-\omega_0)^2}{2\sigma^2}}
\theta(\omega)$
of width $\sigma$ around a frequency $\omega_0$ is prepared.
Most of the analytic results will not depend on this specific
choice. 

In contrast to the usual models of quantum Darwinism where 
information about the system is uniformly spread among all atomic 
fragments, here it is localized in the time-frequency
plane. Hence, we have to specify how we build our fragments.
To start with, it is always possible to agglomerate atoms by 
randomly choosing them.
This leads to the standard behavior of the (averaged/typical) mutual
information discussed in the usual models of quantum
Darwinism~\cite{Zurek-2009,Riedel-2010,Riedel-2011-1}.
%
\Cref{fig/mutualinformation}-(a) deals with the strong
coupling limit where the wavepackets $\phi^{(s)}$ for $s=0$, $1$ are
well separated ($\sigma\,\Delta\tau>1$) (for the
weak coupling, see \cref{sec/weak-coupling}). 
The plots depict the ratio $I(S,F)/S(S)$ as a function
of the number of atomic fragments agglomerated for various values of the 
parameter $T\sigma$ and of the average photon number $n$. The red curves
correspond to coherent radiation whereas the blue curves correspond to
Fock states. At low photon number ($n\lesssim 2$), no Darwinian plateau
appears, an expected result since the low photon number prevents
broadcasting of information about $S$ among many fragments. At large
photon number ($n\gtrsim 8$), we see that a Darwinian plateau is present
for both the classical and quantum probes, although less prominent for
the latter. In the intermediate regime ($2\lesssim n\lesssim 4$), 
which we call mesoscopic, we see a Darwinian plateau for classical 
radiation and no plateau for quantum radiation. 
%However, since correlation between $S$ and
%atomic fragments is highly heterogenous, the
%Darwinian plateaux may be artificially enhanced by the aggregation of
%atomic fragments carrying very little information about $S$.

\begin{figure*}
	\includegraphics{figures/fig2}
\caption{(Color online) Quantum mutual information $I(S,F)/S(S)$ for both classical and
quantum probes prepared with a Gaussian wavepacket as a function of the 
size $f$ of the fragment at strong coupling ($\sigma\Delta\tau = 6$).
%The various plots corresponds to different average photon numbers
%$n$ and dimensionless size $T\sigma$ of the Shannon atoms of signals.
%
(a) \emph{Random ordering}: 
%Composite fragments are built from a random choice of atomic fragments. 
For $n\leq 2$ no Darwinian plateau is present, independently
of the source whereas it is always present for $n\geq 8$. In the
mesoscopic regime (here $n\simeq 4$), a Darwinian
plateau is present with classical probes but not
with quantum ones.
%
(b) \emph{Correlation ordering}: 
%Fragments are built using the two algorithms described in the Letter.
%The crucial role of correlations 
%in the mesoscopic ($n\simeq 4$) 
%time-resolved ($T\sigma=1/10$) regime is revealed by the algorithmic 
%dependence of the plateau.
In the mesoscopic ($n\simeq 4$) time-resolved ($T\sigma=1/10$)
regime, the buildup of the plateau is algorithmic dependent, revealing the
crucial role of correlations.}
\label{fig/mutualinformation}
\end{figure*}

In order to understand this difference, we have to discuss the role of 
correlations between atomic fragments. In this perspective, we 
introduce two algorithms for agglomerating atomic fragments into 
composite fragments of increasing sizes.
The ``naive'' algorithm orders
the atomic fragments $E_{k,l}$ by decreasing order of their correlation with 
the system~\footnote{In case of equality, a random choice is performed.}: 
$I(S,E_1) \geq \dots \geq I(S,E_n)$. 
Composite fragments are then built sequentially by
aggregating atomic fragments accordingly.
Importantly, this construction only requires the reduced density matrices
$\rho_{SE_k}$ and does not use inter-fragment correlations.
By contrast, the ``smarter'' algorithm fully exploits them.
A composite fragment $F_f$ of size $f$ is recursively built by choosing,
at each step, the atomic fragment such that $F_f$ is most correlated
with~$S$.
Given $F_f = \lbrace E_1,\dots,E_f \rbrace$, $F_{f+1}$ is constructed by
aggregating 
%
\begin{equation}
	E_{f+1} 
=\underset{X}{\mathrm{argmax }}\, I(S,X|E_1,\dots,E_f )
\,
\end{equation}
%
where $I(S,X|E_1,\dots,E_f )$ denotes the conditional mutual
information between $S$ and $X$ \cite{Book-Wilde}. This algorithm 
requires the whole set of reduced density matrices
$\lbrace \rho_{SE_1\dots E_k} \rbrace$. 

For classical probes, both algorithms lead to the same behavior of 
the mutual information. This comes from the absence of
correlations between the various atomic fragments for a given state of
the qubit. The quantum mutual information between $S$ and any 
fragment $F$ of the environment is then equal to 
(see \cref{sec/details/classical})
%
\begin{equation}
I(S,F) =h_2(\bigr\lvert D_{\text{tot}} \bigr\rvert)+
f\left(\bigr\lvert D_F \bigr\rvert\right)-
f\left(\left\vert \frac{D_{\text{tot}}}{D_F}
\right\vert\right)\,,
\label{eq/coherentmutualinfo}
\end{equation}
with the binary entropy function
$
h_2(x) = -\frac{1+x}{2} \log\frac{1+x}{2} - \frac{1-x}{2}
	\log\frac{1-x}{2}
$, 
$\lvert D_{\text{tot}}\rvert = |\langle q\phi^{(0)}(\omega)
| q\phi^{(1)}(\omega)\rangle|$ the total decoherence factor and
$|D_F|= \prod_{(k,l) \in F}|\langle q\phi^{(1)}_{kl}
| q\phi^{(0)}_{kl}\rangle|$ the decoherence factor of the fragment
$F$ which is factorized  
into decoherence factors relative to the atomic fragments
within $F$. Since, in the smarter algorithm, 
we look for the atom of signal $X$
maximizing $I(S,E_1,\dots,E_f X)$ or, equivalently,
minimizing 
$D_{S,E_1,\dots,E_f X} = D_{S,E_1,\dots,E_f} D_X$, this
leads to the same atomic fragment than the
naive algorithm which minimizes $D_X$ at each step.

By contrast, in the case of a quantum probe, the difference between the
$I(S,F)/S(S)$ curves obtained with the two algorithms shows the role 
of inter-fragment correlations. In this case, $I(S,F)$ is obtained 
from the reduced density matrix 
$\rho_{SF} =\sum_{k=0}^n p_F(k) \rho_{SF}^{(k)}$
which is an orthogonal mixture with $p_F(k)$ the probability distribution
 to measure $k$ photons in $F$ and $\rho_{SF}^{(k)}$ the conditional 
state upon measuring $k$ photons (see \cref{sec/details/quantum}). 
Consequently, $I(S,F)$ is an average 
$I(S,F) = \mathbb{E}_{p_F(k)}\big( I_{nk}(S,F) \big)$ of
%
\begin{equation}
I_{n,k}(S,F) =  
	h_2(D_{\text{tot}})
	+ h_2\big( D_{F}(k)\big) 
	- h_2\big(D_{\overline{F}}(n-k)\big)
\,.
\label{eq/Focknk}
\end{equation}
over the photon number probability distribution $p_F(k)$. This quantity
depends on two factors. First the total decoherence 
factor 
$
D_{\text{tot}} = 
	\left(
	\int_0^{+\infty} 
		|\phi(\omega)|^2 \me^{\mi\omega\Delta\tau}
	\,\md\omega
	\right)^n
$
which is the autocorrelation function of the analytic
signal $\phi(t)$ as a function of $\Delta\tau$. Second, the 
``decoherence'' factor relative to the fragment $F$ conditioned to $k$
photons in~$F$, 
%
\begin{equation}
	D_F(k)
	= \sqrt{1-4(1-|g_F|^{2k})p_F(0|k)p_F(1|k)}\,,
\end{equation}
where
$p_F(s) = \frac{1}{2\pi}\int_{\mathbb{R}^2} W_F(t,\omega)W_{\phi}(t-\tau_s,\omega)
\, \md t\md\omega$ is the probe wavefunction $\phi$
translated at time $\tau_s$ filtered by the window $W_F$ associated with
the fragment $F$ (see \cref{sec/annexe-wigner}). For each $s=0,1$, $p_F(s|k)$ is the
\emph{a priori} conditional probability distribution to have $s\in\lbrace 0,1\rbrace$
knowing $k$ photons were measured. It is
a Poisson distribution with
parameter $p_F(s)$. Finally $g_F = a_F/\sqrt{p_F(0)p_F(1)}$
where
$|a_F|=\left| \frac{1}{2\pi}\int_{\mathbb{R}^2} 
	W_F(t,\omega)W_{\phi}(t-\overline{\tau},\omega) \,\me^{\mi\omega\Delta\tau}
\, \md t\,\md\omega\right|
$ with $\bar{\tau}=(\tau_0+\tau_1)/2$
is the interference pattern between both $s=0,1$ contributions
filtered by the fragment $F$. This is analogous to Mandel's first degree
of coherence for the filtered wavefunction over $F$.
%which can be viewed 
%as a measure of the statistical dependence between the two signals 
%$\phi^{(0)}$ and $\phi^{(1)}$.

Let us now comment \cref{fig/mutualinformation}-(b). For classical 
probes, both algorithms give the same results but not for quantum 
probes. First of all, the main features 
observed on \cref{fig/mutualinformation}-(a) also appear
for quantum probes: there is no Darwinian plateau at low photon number
($n\lesssim 2$) and a clear plateau appears at large photon number 
($n\gtrsim 8$). 
%With these two algorithms, the curves do not show the smoothness from
%the randomized choice. This feature comes from the optimization process.
The physically relevant differences appear in the mesoscopic regime.

In this regime, classical probes have a more visible plateau than
the quantum ones and the difference between both algorithms 
for quantum probes is more visible than at lower or higher photon 
number. This is a consequence of the correlation structure which
we can better witness by discussing the role of the reference frame, 
\textit{i.e.} changing the $\sigma T$ parameter.

In the frequency-resolved case
($\sigma T =100$, right column \cref{fig/mutualinformation}-(b)), a 
Darwinian plateau appears at lower photon number
for classical probes than for quantum probes. This reflects the fact
that coherent states record information in the phase $\omega\tau_s$ 
between the different photon-number states within a single mode,
while quantum radiation does not. 
%For a classical wave at fixed $\omega$,
%the information about $s$
%is in the phases $\omega\tau_s$.
It is maximal when the phase difference
$\omega\,\Delta\tau$ is $\pi\pmod{2\pi}$ and minimal for
$0\pmod{2\pi}$. This phase can be recovered by homodyne measurements. 
By contrast, Fock states of a fixed $\omega$ mode are eigenstates 
of the number operator $b^\dagger_\omega b_\omega$ which is a building block of 
the dispersive interaction.
%$\sigma_z \otimes b^\dagger_\omega b_\omega$. 
Consequently, a Fock state of a frequency-resolved mode cannot decohere 
the system and distinguish between the two classical states. The relative 
phase between two fragments can nevertheless be recovered by considering
interference experiments (see \cref{sec/fock/freq-resolved}). 
This explains why both algorithms develop 
a plateau at larger photon number with the quantum probe than with the 
classical probe as well as the difference at low $f$.

In the intermediate regime $\sigma T=5$, we see a clear difference 
between the two algorithms on \cref{fig/mutualinformation}-(b),
with a very long plateau for the naive one for $n\geq 4$. 
%
This slower increase of the mutual information arises because the naive 
algorithm aggregates atomic fragments $E_{k,l}$ according to decreasing
order of $I(S,E_{k,l})$. But these are not necessarily the ones that give a fast
rise of the mutual information $I(S,F)$. The smart algorithm is able 
to find an atomic fragment $E_{k,l}$ giving a larger contribution to $I(S,F)$
even if its atomic mutual information $I(S,E_{k,l})$ is lower than for others
$E_{k',l'}$. Additionally, it is also able to recover quantum information 
more efficiently than the naive one 
(see \cref{sec/Holevo} for a discussion of the Holevo information). This
is different from the typical cases of quantum Darwinism where the 
mutual information, the Holevo information and the
quantum discord have symmetry properties~\cite{Zurek-2013}. These
observations question the meaning of the plateau in realistic contexts 
where the information is localized in the environment. 
Indeed, it is not meaningful to interpret the long plateau 
as a redundancy because it may be artificially extended by aggregating
irrelevant atomic fragments. 
%The time-resolved case will
%clearly highlights the role of correlations for a meaningful reconstruction 
%procedure.
%both for the classical and quantum information about $S$.

In the time-resolved case 
($\sigma T=1/10$, left column of \cref{fig/mutualinformation}-(b)), 
the appearance of the Darwinian plateau at increasing $n$
is expected for both types of probe since, the time delay of the 
outgoing plateau unambiguously determines
the qubit state. Consequently, for increasingly large $n$, more and more
$T\sigma\ll 1$ atoms of signals are able to probe the presence of the 
outgoing pulse. The role of correlations is again
clearly visible: for quantum probes, the ``naive'' algorithm builds a
smaller plateau than with the ``smarter'' algorithm.
%and in both cases, it is smaller than for classical probes. 

Starting from a single atomic fragment, the ``smarter'' algorithm opens 
the time window $F$ around one of the time delays, for instance 
$\tau_0$. Then, the detection of one photon tells us that $s=0$ while
the absence of detection is inconclusive. Consequently, the ``smarter'' 
algorithm builds $F$ to minimize the probability $p_F(0)$ of no 
detection. In this case, $I(S,F)$ is approximately equal to a
classical mutual information associated to a measurement of the 
$\sigma_z$ observable of the qubit and a photon number measurement 
$N_F$ within the time window defined by $F$ (see 
\cref{sec/fock/time-resolved}). Hence, the classical data about~$S$ and
how to access it in~$E$ are fully characterized in this regime. As its 
size $f$ increases, the composite fragment $F$ starts to cover both 
regions around $\tau_0$ and $\tau_1$, leaving an ambiguity on the 
detection time of the photons, a situation observed at lower $f$ with 
the naive ordering algorithm. Then, quantum correlations in the state 
$\rho_{SF}$ are probed and $I(S,F)$ exceeds its
classical bound, ending the Darwinian plateau. 
By exploiting correlations, the ``smarter'' algorithm 
gradually builds a ``path detector'' from photon detection
collecting classical information about $s$
more efficiently than the ``naive'' algorithm.
%: in this regime, 
%a Darwinian plateau is reached at
%lower and ends at larger $f$ than by not
%exploiting correlations.

\emph{Conclusion}
To conclude, we have developed an experimentally relevant model of
a qubit coupled to a transmission line and used time-frequency
analysis to study the resource dependence of quantum Darwinism.
This model, whose predictions may be tested in circuit QED
experiments \cite{Blais-2021-1}, allows to quantitatively analyze the role
of the initial preparation and the choice of reference frame.
%(atoms of signals). 
Three different regimes of quantum Darwinism
are unraveled: a deeply quantum regime where no classical image can
be reconstructed, a robust Darwinian regime present whatever resources
are used and, in between, a mesoscopic regime where the ability to 
reconstruct a classical picture is resource dependent. 
%In the latter case, 
%objectivity strongly depends on the reference frame whereas, in the robust Darwinian
%regime, objectivity is achieved in most of them.
As emphasized in \cite{Giacomini-2019-1,Le-2020-1,Tuziemski-2020-1}, 
this shows that objectivity is indeed tightly linked to the correlation 
structure between the system and its environment. Interestingly, 
with our realistic model, we can identify the scale where
quantum Darwinism is blurred with a physically relevant parameter 
(photon number). Because the mesoscopic regime occurs at low photon 
numbers, this calls for going beyond asymptotic quantum 
Shannon theory in order to fully understand the 
emergence of objectivity in a quantum world when a single
realization is considered and therefore address the quantum
measurement problem \cite{Book-von-Neumann,Schlosshauer-2005-1}.

\begin{acknowledgments}
We acknowledge support from the European Space Agency (Ariadna Study
1912-01) and the ANR project 
QuSig4QuSense (ANR-21-CE45-0012).
A.F. acknowledges support from ERC grant BLACKJACK (ERC-2019-STG-851866) 
and ANR AI chair BACCARAT (ANR-20-CHIA-0002).
We are thankful to Jean-Michel Raimond, Irénée Frérot and Franck
Laloë for their insightful comments.
\end{acknowledgments}

